Explain how to refine the Euler-Maruyama approximation by expanding both the
drift and the diffusion term to the same order.

The Euler-Maruyama was our standard approach to simulate stochastic processes,
meaning given an SDE 
\( dX_t = a(X_t) dt + b(X_t) dW_t,\ X_0 = x_0 \) in \( \mathbb{R}^{d} \)
with solution
\( X_t = x_0 + \int_{0}^{t} a(X_u) du + \int_{0}^t b(X_u) dW_u \)
(so \( a : \mathbb{R}^{d} \to \mathbb{R}^{d},\ b : \mathbb{R}^{d} \to \mathbb{R}^{d\timesm}\)
as \( W \) is assumed to be a \( m \)-dim BM) the scheme discretizes 
\( \hat{X}_{t_{i+1}} = \hat{X}_{t_i} + a(\hat{X}_{t_i})(t_{i+1} - t_i) + b(\hat{X}_{t_i})\sqrt{t_{i+1} - t_i}Z_{i+1}\)
if the distances of the simulation steps \( 0 = t_0 < t_1 < \dots < t_m \) are of fixed distance \( h \) one can replace
\( t_{i+1} - t_i = h \).

One can optimize the precision even further via Higher-Order expansion from the Taylor

series however this looks slightly different as usual since we're dealing with Stochastic calculus
according to Ito.

\( d(b(X_t)) = \mu_b(X_t) + \sigma_b(X_t)  + b'(X_t)b(X_t) dW_t \)
where \( \mu_b(X_t) = b'(X_)a(X_t) + \frac{1}{2}b''(X_t)b^2(X_t)dt \)
and \( \sigma_b(X_t) = b'(X_t)b(X_t) dW_t \)

An Euler approximation \( b(X) \) is given by
\( b(X_u) \approx b(X_t) + \mu_b(X_t)(u - t) + \sigma_b(X_t)(W_u - W_t) \)
and we have \( \mu_b(X_t)(u - t) \in \mathcal{O}(u - t) \), \(\sigma_b(X_t)(W_u - W_t) \in \mathcal{O}(\sqrt{u - t}) \)
we're only interested in \( \mathcal{O}(h) \) order terms so we can ignore \( \mu_b \).
end up with
\( \int_{t}^{t+h} b(X_u) dW_u = b(X_t)(W_{t+h} - W_t) + b'(X_t)b(X_t)\int_{t}^{t+h} W_u - W_t dW_u \)
and \( \int_{t}^{t+h} W_u - W_t dW_u  = \frac{1}{2}(W_{t+h} - W_t)^2 - \frac{1}{2}h\) according to Ito's Lemma.
put into sampling notation 
\( \hat{X}_{i+1} = \hat{X}_i + a(\hat{X}_i)h + b(\hat{X}_i)\sqrt{h}Z_{i+1} + \frac{1}{2}b'(\hat{X}_i)b(\hat{X}_i)h(Z^2_{i-1} - 1) \)
which is also referred to as the Milstein-Scheme.

What are strong and weak convergence orders?

One can heuristically describe the convergence orders as the following criteria:
Strong criteria
- pathwise proximity of discretized process to original process
Weak criteria
- proximity of the corresponding distributions
for the time interval \( [0, T], h > 0 \) fixed and \( n = \frac{T}{h} \) floored to the nearest integer.

For \( \beta > 0 \)
Strong order of convergence
\( \mathbb{E}[\|\hat{X}_{nh} - X_T\|] \leq ch^\beta \) for constant \(\beta, c, h > 0\) small
Weak order of convergence
\(| \mathbb{E}[f(\hat{X}_{nh})] - \mathbb{E}[f(X_T)]]| \leq ch^\beta \) for constant \( \beta, c > 0\) 
and \( f \mathbb{R}^{d} \to \mathbb{R}) \) a suitable famility of functions
e.g. \( C^{2\beta+2}_p \) functions polynomially bounded \( \exists k, q \forall x \in \mathbb{R}^d : |g(x)| \leq k(1+ \|x\|^q) \)
for order \( 0, 1, \dots, 2\beta + 2 \).

Euler: If the process \( X \) is sufficiently nice, the strong order is \( \frac{1}{2} \) and the weak order is \( 1 \).
Milstein: For a sufficiently regular process, both strong and weak order are \( 1 \).

Explain a simple strategy (extrapolation) to improve the weak convergence order.
Does the discretization error depend on the simulation methodology? Where does
the simulation methodology matter?

Idea
Combine two 1st-order schemes at different levels of discretization to produce a 2nd-order scheme

Notation
\( X^h \) is the discretized process with step size \( h \)
\(X^h(t)\) is the value of process at time \( t \)

Fact
If both the process \( X \) and the function \( f \) are sufficiently regular, then
\( \mathbb{E}[f(X^h(T))] = \mathbb{E}[f(X(T))] + ch + o(h) \)

This implies:
\( \mathbb{E}[f(X(T))] = 2\mathbb{E} [f(X^h(T))] - \mathbb{E}[f(X^{2h}(T))] + o(h) \quad (*) \)
\( 2\mathbb{E}[f(X^h(T))] - \mathbb{E}[f(X^{2h}(T))] \) is the object that needs to be computed from MC in order to reduce discretization error.

(i) Reduction of discretization error in (*)  involves difference of two expectations; the discretization error does not
depend on the way how the two expectations are simulated - independently or with dependence.
ii) However, the Monte-Carlo error can be reduced if the simulation of both expectations involves suitable
dependence:
\( \mathbb{E}[f(X^h(T))] - \mathbb{E}[f(X^{2h}(T))] = \mathbb{E}[2f(X^h(T)) - f(X^{2h}(T))] \)
\( \text{Var} [2f(X^h(T)) - f(X^{2h}(T))] \)
\( = \mathbb{E}[(2f(X^h(T)) - f(X^{2h}(T)))^2] - (\mathbb{E}[2f(X^h(T)) - f(X^{2h}(T))])^2 \)
\( = 4 \text{Var}[f(X^h(T))] + \text{Var}[f(X^{2h}(T))] - 4 \text{Cov} (f(X^h(T)), f(X^{2h}(T))) \)
\( \Rightarrow \) Positive dependence is useful.

Possible strategy:
\( X^h \): Brownian increments \( \sqrt{h} Z_1, \sqrt{h} Z_2, \sqrt{h} Z_3, ... \) with \(Z_1, Z_2, ... \sim N(0,1)\) i.i.d
\( X^{2h} \) Brownian increments \( \sqrt{h}(Z_1 + Z_2), \sqrt{h} (Z_3 + Z_4), ... \)

More generally
Given a probability space \( (\Omega, \mathbb{F}, \mathbb{P} \), a filtration is a family \(\{\mathcal{F}_t\}_{t \in T}\) of sub-σ-algebras of \( \mathbb{F} \) indexed by a totally ordered set \( T \) (typically representing time), such that it satisfies the following properties:
Increasing: For any \( s, t \in T \) with \( s \leq t \), we have \( \mathcal{F}_s \subseteq \mathbb{F}_t \)
Adaptedness: \( \mathcal{F}_{0} \) contains all the null sets of \( \mathbb{F} \) (Optional, but often assumed).
Right-continuity: (Optional, but often assumed) For any \( t \in T \), we have \( \mathcal{F}_t = \cap_{s > t} \mathcal{F}_s \)

Theorem (Arbitrage-Free Prices): The set \( \mathcal{C}_{0} \) of all arbitrage-free prices  \( C_0 \) of a contingent claim \( C_T \) at time \( 0 \) is non-empty and given by  \( \mathcal{C_0} = \{S^{0}_{0}\mathbb{E}_{\mathbb{Q}}[\frac{C_{t}}{S^{0}_{T}}] | \mathbb{Q} \in \mathcal{Q}\} \) such that \( \mathbb{E}_{\mathbb{Q}}[\tilde{C}_T] < \infty\) 


with bounds \( \inf C_0 = \inf_{\mathbb{Q} \in \mathcal{Q}} S_0^0 \mathbb{E}_{\mathbb{Q}}[\frac{C_T}{S^0_T}] \)
and similarly \( \sup C_0 = \sup_{\mathbb{Q} \in \mathcal{Q}} S_0^0 \mathbb{E}_{\mathbb{Q}}[\frac{C_T}{S^0_T}] \)

Meaning for the claim \( C_t \) exists a real number \( r \in \mathbb{R} \) and portfolio/trading strategy \( \nu(t) \) s.t.
\( C_T = r + \int_{0}^{T} p(t) dX_t \), where \( X_t \) encodes the stochastic processes of the assets that make up the claim
(e.g. \( (X_t)_1 \coloneqq  \) saving account \( (X_t)_2 \coloneqq \) some stock.
each 

Convergence in distribution \( d \)
If \( X_n \sim F_n \) then this means
\( \lim_{n \to \infty} F_n(x) = F(x) \)
for \( X \sim F \)
Converge in probability 
\( \lim_{n \to \infty} \mathbb{P}[|\hat{X}_n - X| > \varepsilon] = 0 \)
Almost sure convergence 
\( \mathbb{P}[\lim_{n \to \infty} \hat{X}_n = X] = 1\)
Note that the position of the limes matters a lot.
Converge in \( r \)-th mean (\(\mathcal{L}^r\))
\( \lim_{n \to \infty} \mathbb{E}[|\hat{X}_n - X|^r] = 0\)
Don't forget \( p \implies a.s. \implies d \)
and \( L^r \implies L^s \implies a.s. \) for \( r \geq s \)


Let \(\bar{X}_n = \sum_{i=1}^n X_i\) denote the sample mean for i.i.d. draws from a r.v. \(X\), the weak law of large numbers
states \(\bar{X}_n \xrightarrow{p} E[X]\) for \(n \to \infty\)

The Lindeberg-Levy CLT sates for the same setup
\( \sqrt{n}(\bar{X}_n - E[X]) \xrightarrow{d} \mathcal{N}(0, \sigma^2) \).

Moreover the Markov/Chebyshev inequality states for a measure space \((\Omega, \mathcal{F}, \mu)\), measurable function \( f \)
and \( 0 < p < \infty \)
\( \mu(\{\omega \in \Omega : |f(\omega)| \geq \varepsilon\}) \leq \frac{1}{\varepsilon^p}\int_\Omega |f|^p d\mu \).
This can be generalized for a non-negative and non-decreasing \( g(t) \neq 0 \) to 
\( \mu(\{\omega \in \Omega : f(\omega) \geq \varepsilon\}) \leq \frac{1}{g(\varepsilon)}\int_\Omega g \circ f d\mu \).

Since 


Girsanovs theorem application in finance

This matters in quantitative finance, since it allows us to turn any 
asset whose diffusion is given by some drifting BM (assumption of the Black-Scholes model) 
into a martingale so the fundamental theorem of asset pricing applies.
In particular let \( M_t \) be the risk free asset \( dM_t = r_t M_t dt \)
and the asset of interest \( dS_t = S_t(\mu_t dt + \sigma_t dW_t \)
then \( \frac{S_t}{M_t} = X_t\)
with \( dX_t = X_t((\mu_t - r_t)dt + \sigma_t dW_t \)
factoring terms \( dX_t = \sigma_t X_t(\frac{\mu_t - r_t}{\sigma_t}dt + dW_t) \)
then Girsanov allows us to find a suitable probability measure \( \mathbb{Q} \)
equivalent to the original \( \mathbb{P} \) under which
\( \tilde{W}_t = W_t + \int_{0}^{t}\frac{\mu_t - r_t}{\sigma_t} dt \)
is a driftless BM and substituting this into our \( X_t \) gives
\( dX_t = \sigma_t X_t(\frac{\mu_t - r_t}{\sigma_t}dt + d\tilde{W}_t - \frac{\mu_t - r_t}{\sigma_t}dt) = \sigma_t X_t d\tilde{W}_t \)


Explain the notation used in the lecture for transition kernels of Markov processes.

Let \( (S,\mathcal{S}) \) be the state space and \( \{X_t\}_{t \in T}\) the markov process.
For \( B \subset \mathcal{S} \), and \( s \leq t \), a transition kernel \( \mu_{s,t} \)
denotes
\( \mu_{s,t}(X_s, B) = \mathbb{P}[X_t \in B|\mathcal{F}_s] = \mathbb{P}[X_t \in B|X_s] \) a.s.
One dimensional distributions are written as

\( \nu_t = \mathcal{L}(X_t) \) and \( \nu_0 = \mathcal{L}(X_t) \) is the initial distribution of the Markov process.

Consider kernels \(\mu, \nu\) from \(S\) to \(S\) 
► Kernel from \(S\) to \(S^2\)
\(\forall B \in S^2: (\mu \otimes \nu)(s, B) := \int \mu(s, dt) \int \nu(t, du) 1_B(t, u)\)

This can, of course, iteratively be generalized to kernels from \(S\) to \(S^n\)
► Kernel from \(S\) to \(S\)
\(\forall B \in S: (\mu \nu)(s, B) := (\mu \nu)(s, S \times B) := \int \mu(s, dt) \nu(t, B)\)

For \( A \in \mathcal{S}^{T} \) (maps from \( T \) to the \( \sigma \)-algebra \( \mathcal{S} \)) and \( x \in S \)
any initial distribution \( \nu \) the law of the Markov process with fixed transition semigroup
can now be represented as a mixture

\( \mathbb{P}_{\nu}A = \int_{\mathcal{S}}\mathbb{P}_{x}(A)\nu(dx) \)

Put a little bit more concrete this states if we want to evaluate the probability
of the Markov chain to be in the set \( A \), (which consists of potential paths of the markov chain) 
this is the same as measuring the probability of going from \( x \) into a certain chain in \( A \)
and weighing it according to the probability of starting in \( x \) (which is given by \( \nu \)).


Explain the Cox-Ross-Rubinstein model and its properties

which is also better known as the Binomial options pricing model.

We start from the initial price \( S_0 \) and fix \( u \geq 1,\ 0 \geq d \leq 1 \) the percentage 
the asset might go up or down respectively. We repeat this for a number of fixed time steps
and visualize the movements in a tree, where the nodes are hierarchically ordered by time.
The final nodes contain the potential outcomes of \( S_T \).

Note that multiple ways might lead to the same endresult at a time step \( t \)
e.g. (u,d), (d,u) so we fuse them into one block for simplicity.

We have the no-arbitrage condition \( d < 1 + r < u \), (\( r \) is the usual risk-free interest rate)
and we can calculate the risk neutral probabilities in one step via the formula \( p = \frac{1+r - d}{u-d} \)
for the price going up and \( 1 - p \) for going down.


Give some identities for the Variance, Covariance and 
define the conditional expectation.

\( \text{Var}[X] = \mathbb{E}[(X - \mathbb{E}[X])^2] = \mathbb{E}[X^2] - \mathbb{E}[X]^2 \)
\( \text{Var}[X + a] =  \text{Var}[X]\) for constant \( a \)
\( \text{Var}[aX] =  a^2\text{Var}[X]\)
\( \text{Var}[aX \pm bY] =  a^2\text{Var}[X] + b^2\text{Var}[X] \pm 2ab\text{Cov}[X, Y]\)
More generally for any linear combination
\( \text{Var}[\sum_{i=1}^N a_iX_i = \sum_{i,j=1}^N a_ia_j\text{Cov}[X_i, X_j] = \sum_{i=1}^{N} a_i^2\text{Var}[X_i] + 2\sum_{1 \leq i < j \leq N} \text{Cov}[X_i, X_j] \)
The order of the \( i < j \) in the last equality matters.
This simplifies to a great deal if the \( X_i \) are independent, since then
\( \text{Cov}[X, Y] = 0 \)
continuing with the covariance, it satisfies
\( \text{Cov}[X, a] = 0 \), 
\( \text{Cov}[aX, bY] = ab\text{Cov}[X,Y] \)
\( \text{Cov}[X + a, Y + b] = \text{Cov}[X,Y] \)
\( \text{Cov}[aX + bY, cW + dV] = ac\text{Cov}[X,W] + ad\text{Cov}[X,V] + bc\text{Cov}[Y,W] + bd\text{Cov}[Y,V] \)
and its commutative in its arguments.

Let \( (\Omega, \mathcal{A}, \mathbb{P}) \) be underlying probability space, for \( \mathcal{F} \in \mathcal{A} \)
the conditional expectation is defined as the random variable \(\mathbb{E}[X|\mathcal{F}] \coloneqq Y \)
s.t.
1. \( Y \) is \( \mathcal{F} \) measurable
2. For any \( A \in \mathcal{F}\) we have \( \mathbb{E}[X\mathbb{1}_{A}] = \mathbb{E}[Y\mathbb{1}_{A}] \).
Then for any \( B \in \mathcal{A} \), \( \mathbb{P}[B|\mathcal{F}] \coloneqq \mathbb{E}[\mathbb{1}_{B}|\mathcal{F}] \)
is called the conditional probability of \( B \) given \( \mathcal{F} \).
If we write \( \mathbb{E}[X|Y] \) this is understood as \( \mathbb{E}[X|\sigma(Y)] \)

For \( X, Y \) being random variables 
If we have a joint density it holds \( \mathbb{E}[h(X)|Y] = \frac{\int h(x)f(x,Y)\lambda(dx)}{\int f(x,Y)\lambda(dx)} \)
which is equivalent to the simpler (but not general enough) definition via Bayes rule, by noting that \( \int f(x, Y)\lambda(dx) = \mathbb{P}(Y) \)
and setting \( h(x) = \mathbb{1}_{X \in B} \).

A not so obvious rule is the tower property for \( \mathcal{G} \subseteq \mathcal{F} \subseteq \mathcal{A} \)
it holds
\( \mathbb{E}[\mathbb{E}[X|\mathcal{F}]|\mathcal{G}] =  \mathbb{E}[\mathbb{E}[X|\mathcal{G}]|\mathcal{F}] = \mathbb{E}[X|\mathcal{G}] \)
